% Drawing 67/04-1 -- folder 67, batch 04, archive page 70 (PDF page 71).
% Transcription: en haut à gauche, sous les mots écrits en biais : une courbe fermée, un point $b$ sur elle, en haut à droite ; de $b$ part vers l'intérieur un arc épais qui s'arrête en un point $a$, et vers l'extérieur deux branches marquées $\zeta$ et $\bar\zeta$ ; à l'intérieur, en bas, une petite flèche courbe d'orientation.
% Kind: closed curve with a boundary point b, an inner arc to a and two outer branches. Method 2 (schematic redraw from the facsimile). Status: partial.
% Notes: As on the page: the closed curve C as one closed curve through points read on his outline; the point b on it; the arc I from b inwards to its end point a (heavy in its lower half), labelled a; the two short branches leaving b outwards, the left one labelled zeta and the right one zeta-bar; the small curved arrow inside, near the bottom. All four strokes at b start from the same point. Left out: the marginal « NB … » written obliquely across the upper left of the curve, the two blots he made at b, and the title « Lemme 4 ». Still differs: the upper left of the outline is hidden under the oblique writing and is interpolated; the letters C and I are not written on the figure (the transcription's note names the curve and the arc after the text). Redrawn 2026-11-12 (issue 128): to check again.
% Coordinates: 1 unit = 1 mm = 1/100 of the crop width (bbox in page fractions: [0.1, 0.07, 0.38, 0.37]).
% Redrawn by Claude (claude-fable-5-2), issue #128, 2026-11-12; not checked by a human; not compiled (no TeX in the session).
% Rights: reproduces part of the work -- behind RENDER_TRANSCRIPTS, never published by itself.
\documentclass[tikz,border=3pt]{standalone}
\usepackage{amsmath,amssymb}
\usetikzlibrary{patterns,arrows.meta}
\tikzset{pencil/.style={black!45}, struck/.style={black!35}, mpt/.style={circle,draw,inner sep=1.3pt}, pt/.style={circle,fill,inner sep=0.9pt}, opt/.style={circle,draw,fill=white,inner sep=1.1pt}}
\begin{document}
\begin{tikzpicture}[x=1mm,y=1mm,line cap=round,line join=round,font=\scriptsize]
\draw plot[smooth cycle, tension=0.55] coordinates {(17,-90) (20,-72) (30,-55) (42,-47.5) (52,-49.5) (60,-50.5) (67,-54) (76,-62) (82,-80) (82,-100) (78,-118) (70,-131) (58,-138) (45,-140) (33,-135) (24,-125) (18,-108)};
\node[pt] at (67,-54) {};
\node[inner sep=1pt] at (70.5,-48.5) {\small $b$};
\draw plot[smooth, tension=0.55] coordinates {(67,-54) (61,-63) (58,-75)};
\draw[very thick] plot[smooth, tension=0.55] coordinates {(58,-75) (57,-83) (55.5,-91)};
\node[inner sep=1pt] at (61,-91) {\small $a$};
\draw plot[smooth, tension=0.55] coordinates {(67,-54) (62.5,-50) (60,-43)};
\node[inner sep=1pt] at (56.5,-37) {\small $\zeta$};
\draw plot[smooth, tension=0.55] coordinates {(67,-54) (74,-53.5) (79,-47) (80,-42)};
\node[inner sep=1pt] at (82.5,-37) {\small $\bar\zeta$};
\draw[-{Stealth[length=2mm]}] plot[smooth, tension=0.55] coordinates {(44.5,-120) (49,-120.5) (52,-117) (52,-112.5)};
\end{tikzpicture}
\end{document}
